\subsection{Review of basic statistical concepts}

\noindent\textbf{Population}: The set of data (numeric or otherwise) corresponding to the entire collection of units about which information is sought.

\noindent\textbf{Sample}: A subset of the population data that are actually collected in the course of a study.

\subsubsection{Parameters vs Statistics}

\noindent\textbf{Parameter}: is a number that describes a population. Notation usually denoted with Greek letters, e.g. $\mu,\sigma$. A parameter is a fixed number (usually unknown).

\noindent\textbf{Statistic}: is a number that describes a sample. Sample statistics are usually denoted using Roman letters, e.g. $x,s$. A statistic is a variable whose value varies from sample to sample.

\subsubsection{Descriptive statistics – numeric and graphics}

\indent

Many methods are available for summarising data in both numeric and graphical form

Numeric measures:

\begin{itemize}
    \item Measure of location: Mean, Median, Mode for numeric data; Counts, proportions for categorical data.
    \item Measure of spread: Standard deviation, MAD (median absolute deviation), IQR (interquartile range)
    \item Others: Min, Max, Quartile, Five number summaries (used later in boxplot)
\end{itemize}

\subsubsection{Homogeneous vs non-homogenous data types}

\noindent\textbf{Homogeneous}

\begin{itemize}
    \item Vector: Sequence of data elements of the same basic data type
    \item Matrix: Collection of data elements in a 2-dimensional array with rows and columns
\end{itemize}

\noindent\textbf{Non-homogeneous}

\begin{itemize}
    \item List: More general structure containing other objects (including possibly other lists)
    \item Data frame: Used for storing data, each column can be a different basic type. All columns must have the same length.
\end{itemize}



